% ----------------------------------------------------------- Table: steps
\begin{table}[t]
\caption{The seven bench steps, with the agent profile reported as primary and the
evidence grade assigned. Full inclusion and exclusion boundaries for each step are
released with the data.}
\label{tab:steps}
\small
\begin{tabular}{@{}llcc@{}}
\toprule
\textbf{ID} & \textbf{Bench step} & \textbf{Primary} & \textbf{Grade} \\
\midrule
S01 & Specimen reception and accessioning & P1 & E2 \\
S02 & Centrifugation of primary tubes     & P1 & E1 \\
S03 & Cap removal and aliquoting          & P1 & E1 \\
S04 & Automated analyser operation        & P1 & E5 \\
S05 & GeneXpert cartridge loading         & P3 & E1 \\
S06 & MALDI-TOF target preparation        & P3 & E3 \\
S07 & Open-tube and open-plate culture work & P3 & E3 \\
\bottomrule
\end{tabular}
\end{table}

% ----------------------------------------------------------- Table: profiles
\begin{table}[t]
\caption{Agent profiles for dimension B. Scores are properties of the organism that may be
present, not of the step, and are therefore held as scenario parameters. $\SCS$ is the mean
of the three items.}
\label{tab:profiles}
\small
\begin{tabular}{@{}p{0.055\columnwidth}p{0.49\columnwidth}ccc@{}}
\toprule
 & \textbf{Profile} & \textbf{B1} & \textbf{B2} & \textbf{B3} \\
\midrule
P1 & Bloodborne, not inhalational \newline \emph{(HBV, HCV, HIV in routine specimens)} & 1 & 2 & 2 \\
P2 & Droplet/aerosol transmissible, prophylaxis available \newline \emph{(invasive \textit{N.~meningitidis})} & 4 & 2 & 3 \\
P3 & Low-dose inhalational, routinely unsuspected \newline \emph{(\textit{Brucella} spp., \textit{M.~tuberculosis})} & 5 & 4 & 5 \\
\midrule
\multicolumn{2}{@{}l}{$\SCS$:\quad P1 $=1.67$\quad P2 $=3.00$\quad P3 $=4.67$} & & & \\
\bottomrule
\end{tabular}
\end{table}

% ----------------------------------------------------------- Table: MAIN (wide)
\begin{table*}[t]
\caption{Reference scores for the seven bench steps at their primary agent profile.
$\IARS$ is the mean of five release items, with A2 reverse-coded; $\SCS$ the mean of three
consequence items for the profile shown; $\CRS$ the mean of four control items, scored
separately for baseline practice and for the recommended control set. $\RBI$ follows
Eq.~\eqref{eq:rbi} with $\kappa=0.8$. Bands: Low $\le 4$, Moderate $\le 9$, Elevated
$\le 14$, High $\le 19$, Very high $\le 25$. The evidence grade is reported beside each
score and is never an input to it. Conditionality indicates whether the evidence differs
materially between the step performed as intended and the step performed with a named
failure (\textit{f-c}) or technique deviation (\textit{t-c}).}
\label{tab:main}
\small
\begin{tabular}{@{}llcccccccl@{}}
\toprule
& & & & \multicolumn{2}{c}{\textbf{Baseline practice}} & \multicolumn{2}{c}{\textbf{Recommended controls}} & & \\
\cmidrule(lr){5-6}\cmidrule(lr){7-8}
\textbf{ID} & \textbf{Bench step} & $\IARS$ & $\SCS$ & $\CRS$ & $\RBI$ (band) & $\CRS$ & $\RBI$ (band) & $\Delta$ & \textbf{Evidence} \\
\midrule
S07 & Open-tube and open-plate culture work & 4.60 & 4.67 & 1.50 & \textbf{19.60} (Very high) & 3.25 & 13.07 (Elevated) & 6.53 & E3 \\
S06 & MALDI-TOF target preparation          & 4.00 & 4.67 & 1.25 & \textbf{17.73} (High)      & 3.25 & 10.27 (Elevated) & 7.47 & E3, \textit{t-c} \\
S05 & GeneXpert cartridge loading           & 3.60 & 4.67 & 2.75 & 10.27 (Elevated)           & 4.00 & 5.60 (Moderate)  & 4.67 & E1, \textit{t-c} \\
S03 & Cap removal and aliquoting            & 4.00 & 1.67 & 1.25 & 6.33 (Moderate)            & 3.25 & 3.67 (Low)       & 2.67 & E1, \textit{t-c} \\
S01 & Specimen reception and accessioning   & 3.40 & 1.67 & 1.00 & 5.67 (Moderate)            & 2.75 & 3.33 (Low)       & 2.33 & E2, \textit{f-c} \\
S02 & Centrifugation of primary tubes       & 3.40 & 1.67 & 1.75 & 4.67 (Moderate)            & 3.25 & 2.67 (Low)       & 2.00 & E1, \textit{f-c} \\
S04 & Automated analyser operation          & 3.00 & 1.67 & 2.25 & 3.33 (Low)                 & 3.50 & 1.67 (Low)       & 1.67 & \textbf{E5}, \textit{f-c} \\
\bottomrule
\end{tabular}
\end{table*}

% ----------------------------------------------------------- Table: evidence
\begin{table*}[t]
\caption{The evidence base, step by step. Direction indicates whether the finding supports
release or harm (positive), found none (negative), or records an absence of evidence
(gap). Full extracted findings with provenance are released with the data.}
\label{tab:evidence}
\small
\begin{tabular}{@{}llp{0.60\textwidth}l@{}}
\toprule
\textbf{Step} & \textbf{Grade} & \textbf{Basis of the grade} & \textbf{Sources} \\
\midrule
S07 & E3 & 16 probable laboratory-acquired meningococcal cases, 8 fatal; isolate manipulated outside a cabinet in 15 of 16 laboratories; occupational attack rate $\approx$65$\times$ background; open-bench subplating and blood-culture aspiration among affected workers; $>$200 occupational exposures from open-bench manipulation across 10 laboratories. No aerosol measurement located. & \cite{sejvar2005,mmwr2002,mmwr2014,ackelsberg2020,bmbl6} \\
S06 & E3 & Identification platform in use during 4 exposure events accounting for 84\% of personnel exposures; in 3 of 4, the Gram stain suggested streptococci rather than the taught cue; no seroconversions where a cabinet was in use. No aerosol measurement located. & \cite{ackelsberg2020,asmbruc2025} \\
S05 & E1 & Neither sample preparation nor automated cartridge processing produced culturable bioaerosol, against 16 and 325~CFU/m\textsuperscript{3} for acid-fast smear preparation. & \cite{banada2010} \\
S03 & E1 & No organisms recovered by air or surface sampling during opening of spiked blood tubes behind a plexiglass panel; detection limit 4~CFU/m\textsuperscript{3}. & \cite{dubuis2021} \\
S02 & E1 & No organisms recovered during centrifugation of spiked blood, closed-cap tubes in a swing-bucket rotor, detection limit 4~CFU/m\textsuperscript{3}. Centrifuge spills appear among studied accident types; measured rupture intensity 10--100~CFU/m\textsuperscript{3}, below that of drops and injection. & \cite{dubuis2021,ravnholdt2025,liu2023} \\
S01 & E2 & Dropped container produced the highest source intensity of three simulated procedures, $\approx$10\textsuperscript{3}~CFU/m\textsuperscript{3}. Inference required: 30~mL concentrated suspension in a research laboratory, not clinical specimens. & \cite{liu2023,ravnholdt2025} \\
S04 & \textbf{E5} & No study measuring aerosol release from a clinical chemistry or haematology analyser was located. Automated equipment identified in review as a source requiring characterisation. & \cite{ravnholdt2025} \\
\bottomrule
\end{tabular}
\end{table*}
